\chapter{Висновки}

Отже, було досліджено маски хмарності, що надаються сервісом Copernicus Open Access Hub. В якості цільової було використано територію Київської області. У процесі досліджень було виявлено, що ця маска не покриває реальну область хмар і може бути покращена. Для цього було запропоновано використати морфологічний оператор розширення.

У роботі наведено також результати порівняльного аналізу методів, за допомогою яких можна комбінувати вхідні знімки для процедури вилучення хмарних пікселів. З точки зору швидкодії найкращим виявився метод вікном розміром 3, по якості - метод кожний з кожний, який і був використаний для побудови безхмарного композиту для території Київської області.

Використання композитів з покращеною маскою хмарності дозволить підвищити якість розв’язання прикладних задач спостереження Землі.

Також побудована модель відновлення супутникових даних коли вхідні знімки є неповними. Підготовлено навчальні дані та навчено алгоритми машинного навчання. Найкращу якість та точність дає багатошаровий перцептрон, його й рекомендується використовувати як спосіб вирішення даної проблеми.
